\documentclass[11pt,oneside,openany]{book}
\usepackage{zlc_frontend_notes}

\begin{document}
\frontmatter
\zlctitlepage{ZLC FPGA Pulse Streamer 手册}{Artix-7 35T 边沿流送器 / 1-tick 预取 / 无限流式扫描 / JTAG-to-AXI}{RTL、1-tick 预取、双 bank 流式、仿射扫描引擎、模拟总线 DAC、编译上传流程、资源预算}{2026-06-07}
\tableofcontents
\cleardoublepage
\mainmatter

\chapter{写给新读者}

\section{这块硬件是什么}
pulse-streamer 是一个\tfocus{固定 bitstream + 运行时表}的脉冲序列器。烧一次 bitstream 就得到时钟状态机、板上输出引脚、\pyapi{jtag_axi} + \pyapi{axi_bram_ctrl} 数据通路和片上 BRAM；它\tfocus{不}把某个具体实验脉冲写死进 Verilog。之后每次 \pyapi{On Pulse}，主机都把当前的 \pyapi{PulseTableState}/\pyapi{PulseSequence} 编译成一张紧凑的程序映像，经 \tfocus{JTAG-to-AXI} 写进 FPGA 的 BRAM，再用一个 CTRL 寄存器邮箱（COMMAND/STATUS）控制引擎运行。

这是\tfocus{唯一}一套设计：没有变体、没有向后兼容、没有探针控制面、没有逐通道游程引擎、没有加载器 FSM。整条路径就是"主机打包 BRAM 映像 $\rightarrow$ AXI 写入 $\rightarrow$ CTRL 邮箱命令 $\rightarrow$ 边沿表引擎播放"。

\begin{notebox}[目标器件]
本设计面向 \tfocus{Xilinx Artix-7 35T（xc7a35tfgg484-2）}，FGG484 封装。这是一颗 LUT/BRAM 资源有限的小器件，所以所有容量参数都按它来权衡。
\end{notebox}

近似资源类别：

\begin{codeblock}[xc7a35tfgg484-2]
logic cells: 33,280     CLB LUTs: 20,800     CLB FFs: 41,600
BRAM:        1,800 Kb (50 x 36 Kb RAMB36)   DSP: 90
\end{codeblock}

\section{两句话能力总结}
\begin{itemize}
  \item \tfocus{最小脉冲宽度与分辨率 = 1 个 tick（20 ns）}。背靠背的 1-tick 边沿能逐周期、一拍一个地打出来——靠的是一条连续预取流水线把 BRAM 读延迟藏掉（见第三章）。
  \item \tfocus{扫描点数无上限}。常驻一个 2-bank ping-pong 窗口（4096 点），主机在游标后面把空出来的 bank 流式回填下一段（见第五章）。配合 N-slot 仿射扫描，duration / delay / DAC 值可以\tfocus{在一次连续运行里任意组合地无缝扫描}。
\end{itemize}

\section{文件}
\begin{description}
  \item[\filepath{fpga/pulse_streamer/zlc_edge_streamer.v}] 引擎。全局边沿表（三块并行 BRAM：tick 32b / coeff 64b / mask 62b，强制 \pyapi{READ_LATENCY_B=2}）、深度 \pyapi{FIFO_DEPTH}(=\pyapi{RD_LAT}+1=3) 的连续边沿预取、2-bank ping-pong 扫描窗口、仿射有效 tick MAC、模拟总线 DAC 引擎。
  \item[\filepath{fpga/pulse_streamer/zlc_pulse_streamer_top.v}] 顶层。\pyapi{jtag_axi} $\rightarrow$ \pyapi{axi_bram_ctrl} $\rightarrow$ 区段译码的 BRAM（三块并行边沿 + 扫描窗口 + 总线映像）+ CTRL 寄存器邮箱，驱动引擎与板上 62 路输出 + 4x10b DAC 总线。
  \item[\filepath{fpga/pulse_streamer/host/image.py}] 主机$\leftrightarrow$FPGA AXI 写契约 + 容量求解器的\tfocus{单一真相源}。\pyapi{pack_program}/\pyapi{unpack_program} 打包/解码 BRAM 映像，\pyapi{solve_capacity} 从器件 RAMB36/LUT 预算推导几何。RTL localparam 与 create-project Tcl 都从它派生。
  \item[\filepath{fpga/pulse_streamer/host/engine_model.py}] 逐周期行为模型（仓库\tfocus{没有 Verilog 仿真器}，这是硬件前的正确性证明）。
  \item[\filepath{Zou_lab_control/neutral_atom/devices/axi_session.py}] 主机会话 \pyapi{VivadoAxiStreamerSession}：常驻 Vivado \pyapi{hw_axi}，打包上传 + 驱动 CTRL 邮箱 + 流式回填。
\end{description}

\chapter{架构总览}

\section{从主机到引脚的一条链}
真实路径是双电脑结构。控制电脑通过 RPyC 把命令发到 FPGA 电脑上的 server；server 持有一个常驻 Vivado \pyapi{hw_axi} 会话，用 JTAG-to-AXI 把程序映像写进 FPGA 的 BRAM，再写 CTRL 寄存器命令引擎运行。

\begin{figure}[htbp]
  \centering
  \begin{tikzpicture}[
      node distance=5mm,
      box/.style={draw=StarkGray, rounded corners=2pt, fill=StarkPale, align=center, inner sep=4pt, font=\scriptsize, minimum height=9mm, text width=26mm},
      hw/.style={box, fill=StarkTint},
      eng/.style={box, fill=StarkCodeBack, draw=StarkTeal, text width=30mm},
      arr/.style={-{Latex[length=2mm]}, StarkTeal, thick}]
    \node[box] (host) {主机 \pyapi{axi_session}\\pack BRAM image};
    \node[box, right=10mm of host] (jtag) {\pyapi{jtag_axi}\\(hw\_axi)};
    \node[box, right=8mm of jtag] (bram) {\pyapi{axi_bram_ctrl}};
    \node[hw, below=8mm of bram] (regions) {区段译码 BRAM\\TICK / COEFF / MASK\\SCAN(2-bank) / BUS};
    \node[hw, left=8mm of regions] (ctrl) {CTRL 邮箱\\COMMAND / STATUS\\CURSOR / BANK\_*};
    \node[eng, below=8mm of ctrl] (engine) {\pyapi{zlc_edge_streamer}\\预取 + 仿射 MAC\\+ 总线 DAC};
    \node[hw, right=8mm of engine] (out) {62 路 TTL\\+ 4x10b DAC};
    \draw[arr] (host) -- (jtag);
    \draw[arr] (jtag) -- (bram);
    \draw[arr] (bram) -- (regions);
    \draw[arr] (bram) -- (ctrl);
    \draw[arr] (regions) -- (engine);
    \draw[arr] (ctrl) -- (engine);
    \draw[arr] (engine) -- (out);
    \draw[arr, StarkBlue, dashed] (engine.north west) to[bend left=20] node[midway, left, font=\tiny]{CURSOR/STATUS} (ctrl.south west);
  \end{tikzpicture}
  \caption{数据通路：主机经 JTAG-to-AXI 写区段译码的 BRAM 与 CTRL 邮箱；引擎从 BRAM 播放，并把 STATUS/CURSOR 回写邮箱供主机轮询/流式回填。}
  \label{fig:fpga-arch}
\end{figure}

\section{CTRL 寄存器邮箱}
主机不碰任何探针；它只读写一小块 CTRL 寄存器文件（来自 \pyapi{host.image.CtrlWords}）：

\begin{codeblock}[CTRL mailbox]
COMMAND     主机->top   上升沿  LOAD(1) / FIRE(2) / RESET(4) / SAFE(8)
STATUS      top->主机   LOADED(1)/RUNNING(2)/DONE(4)/ERROR(8)/UNDERFLOW(16)
PROG_COUNT             边沿行数
SCAN_COUNT            扫描点总数 N (可超过常驻窗口)
SCAN_ENABLE / REPEAT_FOREVER
LOOP_START / LOOP_COUNT / LOOP_END_TICK / LOOP_END_LO / LOOP_END_HI
BUS_COUNTS / BANK_SIZE / SLOT_COUNT
CURSOR      top->主机   已消费的扫描点数 (驱动流式回填)
BANK_READY  主机->top   bit b = bank b 已装好
BANK0_CHUNK / BANK1_CHUNK 主机->top  各 bank 当前装的是哪一段 chunk
\end{codeblock}

生命周期：\pyapi{prepare}（SAFE、上传映像、arm bank、LOAD）/ \pyapi{fire}（FIRE）/ \pyapi{wait_done}（轮询 STATUS、跟随 CURSOR 流式回填）/ \pyapi{safe_state}（SAFE）。COMMAND 字在每个命令前先清 0，制造干净的上升沿。

\section{一行边沿的含义}
引擎里保存一张全局边沿表，每行三个字段，分别存在三块\tfocus{并行}的 BRAM 里：

\begin{codeblock}[每行边沿]
tick[i]  : 该边沿的基准 FPGA tick (32 bit)            -- TICK BRAM
coeff[i] : NUM_SLOTS 个定点扫描系数 (NUM_SLOTS*16 bit) -- COEFF BRAM
mask[i]  : 该 tick 上完整的输出状态 (62 bit)          -- MASK BRAM
\end{codeblock}

bit 0 对应 \pyapi{channels[0]}。一行的含义是：\tfocus{在这个有效 tick，把所有输出切换成这个 mask}——不是逐 channel 命令。三块 BRAM 并行读，所以一次访问拿到一整行边沿，无需位宽填充；\pyapi{max_edges} 因此可以做大（35T 上 4096 行）。

\section{时钟与对齐}
真实硬件默认 50 MHz，一个 tick 是 20 ns。主机在上传前校验所有 period 持续时间、channel 延时、扫描点都对齐到这个 tick 网格（\pyapi{1e9 / clock_hz}）。32 位 tick 计数器在 50 MHz 下覆盖约 85.9 秒。同步后的 FIRE 路径会引入一个固定的小启动延迟，但那是常数偏移，不是比例因子；脉冲宽度和延时由 tick 计数决定。

\chapter{1-tick 预取流水线}

\section{为什么需要预取}
边沿表放在 block RAM 里，而 block RAM 是\tfocus{同步读}：发出读地址后，数据要 \pyapi{RD_LAT} 个周期后才到。build tcl 用 \pyapi{zlc_force_latency2} 把边沿 BRAM 强制成 \pyapi{READ_LATENCY_B=2}，所以 \pyapi{RD_LAT=2} 是确定的。如果天真地"到点了才去读下一行"，每条边沿之间就至少要空 2 个周期——背靠背的 1-tick 边沿根本打不出来。

\section{解法：深度 FIFO\_DEPTH 的连续预取 + 影子}
引擎维护一个深度 \pyapi{FIFO_DEPTH}(=\pyapi{RD_LAT}+1=3) 的"arm" FIFO，里面装着\tfocus{接下来要打的几条边沿}。它\tfocus{每个周期发一次 BRAM 读}，把读回的边沿压进 FIFO 尾；当 \pyapi{time_count} 追上 FIFO 头那条边沿的有效 tick 时，就在\tfocus{同一周期}把它的 mask 打到输出、并把它弹出。因为 FIFO 深度比读延迟多 1，FIFO 永远不空，背靠背的 1-tick 边沿就能一拍一个地连续命中。

\begin{notebox}[SEED 不变式（行为模型逼出来的关键）]
在每个边界，从"第一条还没输出的边沿"开始\tfocus{播种 \pyapi{FIFO_DEPTH} 条影子}，并且边界当拍\tfocus{不发新读}（占用 == 影子数 $\le$ 深度）。3 条影子刚好够给紧跟着的 1-tick 后继留出时间；2 条影子会少一拍而丢边沿。这条不变式是逐 tick 镜像仿真"逼"出来的。
\end{notebox}

\section{时序图：背靠背 1-tick 边沿一拍一个}
下面把时钟周期画在横轴，信号画在纵轴。边沿表是 4 条 1-tick 边沿（有效 tick = 0,1,2,3，mask = A,B,C,D）。\pyapi{rd@T} 表示"周期 T 发出的 BRAM 读"，它的数据在 \pyapi{T+RD_LAT}（=T+2）落地（\pyapi{<A}=A 落地）。注意 \pyapi{out} 在 tick 0,1,2,3 上恰好逐拍输出 A,B,C,D，中间\tfocus{没有空拍}。

\begin{codeblock}[1-tick prefetch (RD\_LAT=2, FIFO\_DEPTH=3)]
cycle (clk)   :  -2   -1    0    1    2    3    4
time_count    :   -    -    0    1    2    3    4
issue BRAM rd :  rdA  rdB  rdC  rdD   -    -    -      (one read per cycle)
data lands    :   -    -   <A   <B   <C   <D   -       (T+RD_LAT, RD_LAT=2)
arm FIFO depth:   0    1    2    3    2    1    0       (>=1 until drained)
FIFO head edge:   -    -    A    B    C    D    -
out (mask)    :   0    0    A    B    C    D    0
                            ^    ^    ^    ^
                         tick0 tick1 tick2 tick3  <- one 1-tick edge per clock
\end{codeblock}

\begin{itemize}
  \item 周期 $-2,-1$ 是\tfocus{arm 阶段}（reset/LOAD 期间）：连发读把前几条边沿的影子预取进 FIFO，使 \pyapi{time_count} 从 0 开始时 FIFO 已经有货。
  \item tick 0：FIFO 头是 A 且 \pyapi{time_count==tick[A]==0}，\tfocus{当拍}输出 A 并弹出；同一拍 C 的数据正好落地补进 FIFO。
  \item 由于 FIFO 深度 3 $>$ 读延迟 2，弹出与补充永远咬合，tick 1/2/3 连续命中 B/C/D，\tfocus{无空拍}。
\end{itemize}

\begin{notebox}[这条等价是被证明的]
\pyapi{host.engine_model.rtl_mirror_play} 在读延迟 1/2/3 下逐 tick 等于组合参考播放器 \pyapi{reference_play}，并跑了 200 个 fuzz 程序（含背靠背 1-tick、扫描 1-tick、loop 1-tick）。仓库无 Verilog 仿真器，这个逐周期镜像就是硬件前最强证据。
\end{notebox}

\chapter{无缝边界：四个 gapless 重装点}

\section{边界为什么不留缝}
引擎只有\tfocus{一个}全局边沿指针，所以一次 repeat / loop-rewind / scan-advance / start 只是\tfocus{单周期}的指针 + 影子重装。四个 gapless 重装点（start / loop-rewind / scan-advance / repeat）在边界用上一节的 SEED 不变式重新播种 \pyapi{FIFO_DEPTH} 条影子，于是\tfocus{点 k 的最后一条边沿与点 k+1 的第一条边沿在时间上紧挨着，中间没有缝}。

\section{时序图：扫描点 k -> k+1 跨界无缝}
设每个扫描点的边沿模板末尾在相对 tick $L-1$ 输出 mask $Z$，下一点的首条边沿在该点的有效 tick 0（绝对上紧跟在 $Z$ 之后一拍）输出 mask $A'$。引擎在播放点 k 末尾的同时，已经把点 k+1 的首批影子按新 slot 向量重新播种好了：

\begin{codeblock}[scan-advance boundary (gapless)]
abs tick      :  L-3  L-2  L-1 | L    L+1  L+2          (| = scan-advance)
playing point :   k    k    k  | k+1  k+1  k+1
issue BRAM rd : rdY  rdZ  rdA' | rdB' rdC'  -           (k+1 reads cross-issued)
arm FIFO      :  ...   Z   A'  | B'   C'   ...          (reseeded at boundary)
out (mask)    :   X    Y    Z  | A'   B'   C'
                            ^    ^
                       point k  point k+1
                       末edge   首edge  <- 相邻一拍, 无 gap
\end{codeblock}

\begin{itemize}
  \item 边界处（\pyapi{|}）引擎\tfocus{不发新读}，只用已经播种好的影子（A'、B'…）。新 slot 向量在边界当拍生效，所以 A' 的有效 tick 已经按 k+1 的扫描值算好。
  \item loop-rewind（repeat bracket）走的是同一套机制：指针跳回 \pyapi{loop_start_addr}，影子按 loop 起点重新播种，循环接缝同样无缝。
\end{itemize}

\pyapi{host.engine_model} 用 \pyapi{streaming_scan_play} 与 1-tick 边界 fuzz 证明了这一点：跨扫描点 / 跨 loop 的边沿与不分段的参考播放逐 tick 相同。

\chapter{无限流式扫描：2-bank ping-pong}

\section{窗口 + 流式回填}
扫描窗口在硬件里是一个 \tfocus{2-bank ping-pong}（\pyapi{BANK_SIZE}=2048，是 2 的幂；两 bank 共 4096 个常驻扫描点），放在一块 BRAM 里。引擎依次播放点 \pyapi{0..N-1}，用 \pyapi{(idx/BANK_SIZE) mod 2} 选 bank，并把已消费点数写进 \pyapi{CURSOR}。主机读 \pyapi{CURSOR}，把游标已经走过的那个 bank 回填\tfocus{下一段 chunk}，并更新 \pyapi{BANK_READY}/\pyapi{BANK*_CHUNK}。这样扫描点文件虽然有限，但\tfocus{大小无上限}。

\begin{figure}[htbp]
  \centering
  \begin{tikzpicture}[
      bank/.style={draw=StarkGray, rounded corners=2pt, align=center, font=\scriptsize, minimum height=8mm, minimum width=22mm},
      live/.style={bank, fill=StarkTint, draw=StarkTeal},
      refill/.style={bank, fill=StarkCodeBack},
      arr/.style={-{Latex[length=1.8mm]}, thick},
      lbl/.style={font=\tiny, StarkGray}]
    \node[live] (b0) {bank0\\chunk 0};
    \node[refill, right=14mm of b0] (b1) {bank1\\chunk 1};
    \node[lbl, below=1mm of b0] {引擎正在播放};
    \node[lbl, below=1mm of b1] {下一段待播};
    \draw[arr, StarkTeal] ([yshift=2mm]b0.east) -- node[above, lbl]{CURSOR 推进} ([yshift=2mm]b1.west);
    \draw[arr, StarkBlue, dashed] ([yshift=-3mm]b1.east) to[bend right=40] node[below, lbl]{主机回填 chunk 2}([yshift=-3mm]b0.east);
    \node[refill, right=20mm of b1] (note) {游标走出 bank0 后\\bank0 空闲\\回填 chunk 2};
  \end{tikzpicture}
  \caption{2-bank ping-pong：引擎播放当前 bank、推进 CURSOR；主机在游标背后把空出来的 bank 回填下一段 chunk。}
  \label{fig:fpga-pingpong}
\end{figure}

\section{bank\_chunk 握手与晚到 STALL}
引擎只在 \pyapi{BANK_READY} 为真\tfocus{且}该 bank 装的是\tfocus{正确 chunk}（\pyapi{BANK*_CHUNK} 与引擎期望的段号一致）时才推进。如果回填来晚了，引擎\tfocus{停住}（hold 在当前点，置 \pyapi{STATUS_UNDERFLOW}），\tfocus{绝不}播放错误或过期的点；主机回填到位后它再继续。这把"晚到"变成一次安全的 hold，而不是数据损坏。

\begin{codeblock}[streaming refill timeline + late-refill STALL]
CURSOR (pts done):  0 ........ 2047 | 2048 ...... 4095 | 4096 .....
engine plays bank:  bank0          | bank1           | bank0(chunk2)
BANK_READY        :  b0,b1 ready    | b0 freed        | refill in flight
BANK0_CHUNK       :  0              | (host -> 2)     | 2
host refill bank0 :       [.... writing chunk 2 ....]
                                          ^
case A 及时:  chunk2 在游标到 4096 前写好 -> BANK0_CHUNK=2 ready -> 无缝继续
case B 晚到:  游标到 4096 但 bank0 还不是 chunk2
              -> 引擎 STALL (hold 当前点, STATUS_UNDERFLOW=1)
              -> 主机写完 chunk2 + BANK0_CHUNK=2 + BANK_READY
              -> 引擎自动继续, 不丢点/不串点
\end{codeblock}

\section{repeat\_forever 流式重扫}
\pyapi{repeat_forever} 对一个有限（但很大）的流式扫描做反复扫掠：\pyapi{fire} 时启动一个主机\tfocus{后台回填线程}，循环把下一段填进空出的 bank。一次扫掠之内是\tfocus{无缝}的；扫掠与扫掠之间的接缝是一个\tfocus{短暂的安全 hold}（重新从 chunk 0 播种）。\pyapi{host.engine_model.streaming_scan_play} 证明了播放序列、CURSOR 推进、以及晚到 STALL 行为。

\chapter{N-slot 仿射扫描引擎}

\section{核心公式}
每行边沿存一个 base tick 加 \pyapi{NUM_SLOTS} 个定点系数；运行时 FPGA 对每个扫描点计算：

\begin{codeblock}[effective tick]
effective_tick = base_tick + (sum_j coeff_j * slot_j) >>> COEFF_FRAC_BITS
\end{codeblock}

其中 \pyapi{slot_j} 是当前扫描点第 $j$ 个 slot 的值（来自扫描点表那一行），\pyapi{COEFF_FRAC_BITS=8} 是定点小数位数。\tfocus{已经没有 x/y}——slot 是任意多个命名维度 \pyapi{s0,s1,...}，每个绑定一个 duration / delay / DAC 值字段。这同一个 \pyapi{effective_tick} 表达式被边沿 tick、loop 终点 tick、以及总线段起止 tick \tfocus{复用}——这正是为什么 duration、delay、DAC 值能任意组合地同时扫描，且只花一条 MAC 的 LUT。

\section{为什么这是对的取舍}
spec 最初设想的是逐通道游程编码（per-channel）+ 双 bank 全新存储。我们改成把仿射引擎\tfocus{泛化}成 N 个 slot 的流式扫描点表，原因围绕 35T 的现实：
\begin{itemize}
  \item 存储复杂度 O(edges) + O(points $\times$ slots) 已达成 spec 的核心存储目标，但不需要每通道游程的状态机与地址逻辑。
  \item LUT 最省：一条比较器 + 一条仿射 MAC，而不是 62 个通道 player。
  \item 可验证：在无 Verilog 仿真器的前提下，仿射方案能被纯 Python 的逐周期行为模型忠实复刻。
\end{itemize}

\chapter{模拟总线 DAC 引擎}

\section{段表编码}
每条 10-bit DAC 总线有一张小的 LUTRAM 段表（每 tick 组合读），每段编码：

\begin{codeblock}[bus segment]
start_tick / stop_tick : 段的起止 tick (仿射, 带 NUM_SLOTS 个系数, 同 effective_tick)
start_value / stop_value: 段起止 DAC 码 (10 bit)
mode                    : edge/hold 或 ramp
value_select (dual)     : start 与 stop 各一个选择子
\end{codeblock}

\pyapi{edge/hold} 段把输出保持在某个值；\pyapi{ramp} 段在 start 到 stop 之间\tfocus{线性插值}，引擎本地每步走 1 个码（不额外占边沿行）。\tfocus{dual value\_select} 让 start 与 stop 各自可以指向一个扫描 slot——于是一条 ramp 能\tfocus{同时扫两个端点}（扫描-A $\rightarrow$ 扫描-B），一条 edge/hold 段能跟随一个扫描的 DAC 码。

\section{时序图：ramp 插值与扫描端点}
横轴是 tick，纵轴是 DAC 码。下图画一条从 \pyapi{start_value} 到 \pyapi{stop_value} 的 ramp；当端点被绑定到 slot 时，不同扫描点画出不同斜率/高度的同一条段（无需新增边沿行）。

\begin{figure}[htbp]
  \centering
  \begin{tikzpicture}[font=\scriptsize, >=Latex]
    \draw[->, StarkGray] (0,0) -- (7.2,0) node[right]{tick};
    \draw[->, StarkGray] (0,0) -- (0,3.3) node[above]{DAC code};
    % axis ticks
    \draw[StarkGray] (1,0.05) -- (1,-0.05) node[below]{\tiny start};
    \draw[StarkGray] (5,0.05) -- (5,-0.05) node[below]{\tiny stop};
    % point p0 ramp: start_value low -> stop_value high
    \draw[StarkTeal, very thick] (1,0.6) -- (5,2.0);
    \node[StarkTeal, font=\tiny, above left] at (5,2.0) {p0: A0 -> B0};
    \fill[StarkTeal] (1,0.6) circle (1.3pt);
    \fill[StarkTeal] (5,2.0) circle (1.3pt);
    % hold before/after
    \draw[StarkTeal, thick, dashed] (0.2,0.6) -- (1,0.6);
    \draw[StarkTeal, thick, dashed] (5,2.0) -- (6.6,2.0);
    % point p1 ramp (scanned both endpoints): higher start, higher stop
    \draw[StarkBlue, very thick] (1,1.2) -- (5,3.0);
    \node[StarkBlue, font=\tiny, above left] at (5,3.0) {p1: A1 -> B1 (scanned)};
    \fill[StarkBlue] (1,1.2) circle (1.3pt);
    \fill[StarkBlue] (5,3.0) circle (1.3pt);
    % per-step note
    \node[StarkGray, font=\tiny, align=center] at (3,0.25) {引擎每步 +1 码\\(本地插值, 不占边沿行)};
  \end{tikzpicture}
  \caption{ramp 段：引擎在 start/stop tick 之间本地线性插值。端点绑定 slot 时，不同扫描点（p0/p1）画出不同起止——dual value\_select 让一条段同时扫两个端点。}
  \label{fig:fpga-ramp}
\end{figure}

\pyapi{host.engine_model.bus_play} 逐 tick 证明了总线/ramp 引擎，包含被扫描的 ramp 端点。

\chapter{容量与资源预算}

\section{35T 上解出的容量}
容量由 \pyapi{host.image.solve_capacity(part, channel_count)} 求解（单一真相源），目标 \tfocus{$\le$90\%} RAMB36。35T 上解出：

\begin{codeblock}[solve\_capacity('xc7a35tfgg484-2')]
4096 edges + bank_size 2048 (4096 resident) + UNBOUNDED streaming
RAMB36: 39/50 = 78%      LUT: 26%   FF: 12%   DSP: 9%
\end{codeblock}

三块并行边沿 BRAM（tick/coeff/mask）+ 扫描窗口 BRAM 占了 RAMB36 预算的大头；\tfocus{总线段表是 LUTRAM（distributed），不占 RAMB36}——因为总线/ramp 引擎每 tick 组合读它们，搬进 BRAM 会破坏组合读。Vivado \pyapi{report_utilization} / \pyapi{report_timing_summary} 是最终权威，Python 估算只是预算指引。

\chapter{编译与上传流程}

\section{build / program / run}
FPGA 侧只有一个入口 \filepath{fpga\textbackslash build\_and\_program.bat}（内部跑 \filepath{create_project.tcl} $\rightarrow$ \filepath{program_fpga.tcl}），server 入口是 \filepath{fpga\textbackslash run\_server.bat}（jtag-axi 后端）：

\begin{codeblock}[PowerShell]
.\fpga\build_and_program.bat --check    # 不连板的 HDL 综合 + 容量自查
.\fpga\build_and_program.bat            # build + program
.\fpga\run_server.bat --check-config    # 打印 project/bit/ltx/xdc/channels/clock/capacity
.\fpga\run_server.bat                    # 启动常驻 hw_axi server (host 0.0.0.0, port 18861)
\end{codeblock}

\pyapi{create_project.tcl} 生成 \pyapi{jtag_axi} + \pyapi{axi_bram_ctrl} + 5 块 BRAM，并用 \pyapi{zlc_force_latency2} 把边沿 BRAM 强制成 \pyapi{READ_LATENCY_B=2}。产物 \filepath{zlc_pulse_streamer_top.\{bit,ltx\}} 落在 \filepath{fpga\textbackslash build\textbackslash pulse_streamer\textbackslash pulse_streamer.runs\textbackslash impl_1}。server 加载 \filepath{.ltx} 是为了在比特流里定位 \pyapi{jtag_axi} 核（即 \pyapi{hw_axi}），所以 bit 与 ltx 必须来自\tfocus{同一次编译}。

\section{一次上传里发生了什么}
\begin{enumerate}
  \item 主机校验程序：\pyapi{validate_pulse_streamer_program} 检查边沿/扫描点/位宽/通道不越界，并校验\tfocus{全局有效 tick 单调}（扫描不会把合并后的边沿顺序打乱；会打乱顺序的扫描被拒绝，而不是悄悄丢点）。
  \item \pyapi{host.image.pack_program} 把整张程序打包成 32-bit BRAM 映像：CTRL 标量 + 三块并行边沿表 + 2-bank 扫描窗口 + 总线段映像。
  \item \pyapi{prepare}：发 SAFE $\rightarrow$ AXI 写映像 $\rightarrow$ arm 两个扫描 bank（写 \pyapi{BANK_READY}/\pyapi{BANK*_CHUNK}）$\rightarrow$ 发 LOAD（等 \pyapi{STATUS_LOADED}）。
  \item \pyapi{fire}：发 FIRE；对流式/重复扫描，启动后台回填线程。
  \item \pyapi{wait_done}：轮询 STATUS（DONE 位），并跟随 CURSOR 回填空出来的 bank。\pyapi{repeat_forever} 永不置 DONE，所以必须给 timeout 或主动 stop。
\end{enumerate}

\begin{notebox}[全部在硬件前被 Python 验证]
\pyapi{pack_program}$\leftrightarrow$\pyapi{unpack_program} 往返、\pyapi{engine_model} 的逐周期等价（1-tick 预取 / 边界无缝 / 流式 ping-pong / 总线 ramp）、以及 \pyapi{VivadoAxiStreamerSession} 的 pack$\rightarrow$upload$\rightarrow$LOAD$\rightarrow$fire 流程（Tcl 执行器可注入，无需 Vivado）都有契约测试覆盖。仓库没有 Verilog 仿真器，这套 Python 证据是硬件 bring-up 前的正确性保证。
\end{notebox}


\end{document}
